\section{Conclus\~ao e Trabalhos Futuros}
\label{sec:conclusao}

Este trabalho propôs um pipeline que liga as opiniões publicadas sobre audiências públicas à fala de
quem as emitiu, condensa a fala de cada participante em um perfil escrito somente a partir dela e usa
esse perfil para simular como a pessoa se posicionaria, com a base da estimativa declarada e uma
justificação conferida. A UDV fornece evidência localizada para 2.105 das 2.203 opiniões, mas, mesmo
quando a opinião contém a citação, a sentença de maior cosseno coincide com a passagem citada em
61,2\% dos casos, e a precisão da camada semântica depende da validação humana. <Síntese dos
resultados da simulação.>

A sustentação das falas simuladas precisa ser avaliada por leitura humana, pois a conferência
automática cobre a existência dos itens citados e a literalidade das cópias, mas não o sentido. Ficam
também em aberto o efeito dos exemplos de fala, que a múltipla escolha não mede, e a comparação entre
perfis gerados por modelos e prompts diferentes. O mesmo protocolo pode ser aplicado a outras fontes
de fala pública dos parlamentares, desde que o gabarito continue vindo de audiências posteriores às
usadas nos perfis.
